\paragraph{}In this lab we run transformer open and short circuit tests in order to
produce its cantilever equivalent. In an ideal transformer there are no losses in
the system. In practice however there are losses which come from 4 placed. Inside the
transformer core there are losses from hysteresis ($R_c$), which is heat lost due to the
alternating magnetic field polarizing the iron within it, and from eddy currents ($X_m$)
produced by the magnetic field inside the iron, which also dissipate energy as heat.
Additionally there are losses from the transformers copper winding ($R_{eq}$) and
flux leakage, where some of the magnetic flex escapes and is left unused. The cantilever
equivalent is a model which tells us how the transformer behaves in the real world, and
includes these 4 kinds of losses. While performing the open and short circuit tests we
will record the primary and secondary voltages, current through in the primary side,
as well as the real and reactive power. These measurements allow us to calculate the
resistor and inductor values in the cantilever equivalent. Providing us with a realistic
model of the transformer.
